% Схема алгоритма вывода одного узла плана. Вызывается рекурсивно для каждого
% ингредиента рецепта и для недостающего инструмента.
\begin{tikzpicture}[
  font=\footnotesize,
  node distance=5mm,
  term/.style={draw, rounded rectangle, align=center, minimum width=3.4cm, minimum height=8mm},
  proc/.style={draw, align=center, minimum width=6.2cm, minimum height=9mm, inner sep=3pt},
  % Предопределённый процесс: прямоугольник с двумя внутренними вертикальными линиями.
  sub/.style={proc, inner xsep=4mm, append after command={
    \pgfextra{
      \draw ([xshift=2mm]\tikzlastnode.north west) -- ([xshift=2mm]\tikzlastnode.south west);
      \draw ([xshift=-2mm]\tikzlastnode.north east) -- ([xshift=-2mm]\tikzlastnode.south east);
    }}},
  cond/.style={draw, diamond, aspect=3.2, align=center, inner sep=0pt, minimum width=6.2cm},
  exit/.style={draw, align=center, text width=3.4cm, minimum height=9mm, inner sep=3pt},
  arr/.style={-{Stealth[length=1.8mm]}, semithick},
  ylbl/.style={font=\scriptsize, above, pos=0.25},
  nlbl/.style={font=\scriptsize, right, pos=0.4},
]

% --- основная ветвь ---
\node[term] (start) {Начало: предмет $X$, количество $N$,\\запасы, стек вывода, глубина};
\node[proc, below=of start] (take) {Взять из запасов $\min(\text{запас}_X, N)$,\\остаток потребности $N'$};
\node[cond, below=of take] (done) {$N' = 0$?};
\node[cond, below=of done] (limit) {Превышена глубина\\или бюджет перебора?};
\node[cond, below=of limit] (cycle) {$X$ уже в стеке вывода?};
\node[sub, below=of cycle] (recipes) {Для каждого рецепта $X$: вывод ингредиентов\\на копии запасов (рекурсивно)};
\node[sub, below=of recipes] (sources) {Для каждого источника добычи $X$: достроить\\недостающий инструмент (рекурсивно)};
\node[cond, below=of sources] (any) {Найден хотя бы\\один способ?};
\node[proc, below=of any] (best) {Выбрать лучший вариант по выполнимости,\\числу тупиков и трудоёмкости; записать развилку};
\node[proc, below=of best] (apply) {Принять запасы выбранного варианта};
\node[term, below=of apply] (end) {Конец: шаг плана};

% --- исходы, завершающие вывод досрочно ---
\node[exit, right=12mm of done] (have) {Шаг <<есть>>};
\node[exit, right=12mm of limit] (lost) {Шаг <<не получено>>:\\предел вывода};
\node[exit, right=12mm of cycle] (loop) {Шаг <<не получено>>:\\цикл};
\node[exit, right=12mm of any] (none) {Шаг <<не получено>>:\\нет способа};

% --- переходы основной ветви ---
\foreach \a/\b in {start/take, take/done, recipes/sources, sources/any, best/apply, apply/end}
  \draw[arr] (\a) -- (\b);
\draw[arr] (done) -- node[nlbl] {нет} (limit);
\draw[arr] (limit) -- node[nlbl] {нет} (cycle);
\draw[arr] (cycle) -- node[nlbl] {нет} (recipes);
\draw[arr] (any) -- node[nlbl] {да} (best);

% --- досрочные исходы ---
\draw[arr] (done) -- node[ylbl] {да} (have);
\draw[arr] (limit) -- node[ylbl] {да} (lost);
\draw[arr] (cycle) -- node[ylbl] {да} (loop);
\draw[arr] (any) -- node[ylbl] {нет} (none);

% Все досрочные исходы сходятся к концу алгоритма по общей линии справа.
\coordinate (rail) at ($(have.east) + (5mm, 0)$);
\foreach \n in {have, lost, loop, none}
  \draw[semithick] (\n.east) -- (\n.east -| rail);
\draw[arr] (rail) |- (end.east);

\end{tikzpicture}
